%% CAS-style main document reconstructed from the original IEEE project.
%% Set this file as the main document in Overleaf.
\documentclass[a4paper,fleqn]{cas-dc}

%% Reference-project style packages, with the original manuscript's algorithm syntax retained.
\usepackage[authoryear,longnamesfirst]{natbib}
\usepackage{amsmath,amsfonts,amssymb}
\usepackage{amsthm}
\usepackage{graphicx}
\usepackage{booktabs}
\usepackage{tikz}
\usetikzlibrary{positioning,arrows.meta,fit,backgrounds}
\usepackage{algorithm}
\usepackage{algorithmic}
\usepackage{float}
\usepackage{xcolor}
\usepackage{placeins}
\usepackage{balance}
\hypersetup{colorlinks=true,citecolor=blue,linkcolor=blue,urlcolor=blue}
\usepackage{cleveref}

\newtheorem{lemma}{Lemma}
\newtheorem{theorem}{Theorem}
\newtheorem{observation}{Observation}
\newtheorem{definition}{Definition}
\newcommand{\eg}{\it e.g., }
\newcommand{\etal}{\it et~al. }
\newcommand{\ie}{\it i.e., }

\begin{document}
\let\WriteBookmarks\relax
\def\floatpagepagefraction{1}
\def\textpagefraction{.001}

% Short front-matter fields, matching the reference project structure.
\shorttitle{Data-aware DAG scheduling across the Compute Continuum}
\shortauthors{Gupta et~al.}
\title[mode = title]{Data-Aware Scheduling of Approximate DAG Workflows across the Compute Continuum}

% Author 1: corresponding author
\author[1]{aa}
\cormark[1]
\ead{[aa.email@university.edu]}

% Author 2
\author[1]{[Second Author Full Name]}
\ead{[second.author@university.edu]}

% Author 3, if applicable
\author[2]{[Third Author Full Name]}
\ead{[third.author@institution.edu]}

% Affiliation 1
\affiliation[1]{
  organization={[Department or School], [University Name]},
  city={[City]},
  state={[State/Province]},
  country={[Country]}
}

%% TODO(author): supply the real organisation for affiliation 2, or renumber the
%% third author to [1] and delete this block. Placeholder kept so the file compiles.
\affiliation[2]{
  organization={[TODO: Department], [TODO: Institution]},
  city={[TODO: City]},
  country={[TODO: Country]}
}
\begin{abstract}
The Compute Continuum connects edge, fog, cloud, and high-performance computing (HPC) resources so that scientific workflows can place tasks according to latency, data locality, resource capability, and resilience requirements. This creates a scheduling challenge for data-intensive directed acyclic graph (DAG) workflows: a placement that is computationally attractive can still violate a deadline when intermediate data traverse bandwidth-constrained links or when a resource fails. This paper extends Gayani Gupta's IEEE work on federated-fog DAG allocation by positioning federated fog as the latency-sensitive tier of a Compute-Continuum-aware eScience workflow.

We present a data- and bandwidth-aware approximate-computing framework that models task placement, execution fraction, precedence constraints, communication delay, resource eligibility, and workflow deadlines. The framework uses a mixed-integer linear programming (MILP) benchmark and retains scalable randomized-relaxation, greedy, and genetic-algorithm schedulers. The model explicitly incorporates task-level data size and link-level bandwidth, enabling an orchestrator to distinguish local and remote placements according to both compute and communication cost.

The supplied simulation results compare the approaches under four conditions: slack-factor variation, federated resource scale, workflow scale, and workflow-node failures. The revised methodology also specifies FAIR-oriented workflow metadata, configuration capture, provenance records, and reproducibility artifacts without claiming release of resources not supplied with the study. Together, the formulation, algorithms, proofs, and evidence provide a foundation for interoperable, resilient, and resource-aware scientific workflow execution across the Compute Continuum.
\end{abstract}

\begin{highlights}
\item Extends Gayani Gupta's IEEE federated-fog workflow paper to a Compute-Continuum-aware eScience setting.
\item Jointly models task placement, approximate execution, data-transfer size, bandwidth, eligibility, and deadlines.
\item Retains exact, randomized, greedy, and genetic scheduling methods with original proof material in an appendix.
\item Compares deadline flexibility, resource scale, workflow scale, and failure resilience using the supplied simulation results.
\item Specifies FAIR-oriented metadata, provenance, and reproducibility artifacts for workflow reuse.
\end{highlights}

\begin{keywords}
Compute Continuum \sep High-Performance Computing \sep Scientific Workflows \sep Federated Fog \sep Directed Acyclic Graphs \sep Approximate Computing \sep Resource-Aware Scheduling \sep FAIR Workflows \sep Provenance
\end{keywords}

\maketitle

\input{Sources/intro}
\input{Sources/relatedwork}
\input{Sources/systemmodel}
\input{Sources/reproducibility}
\input{Sources/orchestration}
\input{Sources/resourceAllocation}
\input{Sources/results}
\input{Sources/validity}
\input{Sources/conclsn}

\appendix
\section{Original Formulation Proofs and Linearization}
\label{app:original_proofs}
\input{Sources/appendix_proofs}

\section*{Acknowledgments}
We would like to thank the anonymous reviewers for their valuable feedback and insightful suggestions. Their comments have helped improve the quality and clarity of this work.

This research is supported by the National Science Foundation (NSF).

\printcredits
\balance
\bibliographystyle{cas-model2-names}
\bibliography{cas-refs}

\end{document}
